%% kapa: exam=1 round=20 year=2009 period=1 form=A subject=civil_law subject_ko=민법 seq=33 number=33
%% source: Q-Net 감정평가사 기출문제 (cst00309) artlSeq=5024178 file=extracted/1008/5024178/2009년 제20회 감정평가사 1차시험 기출문제/1교시.hwp
%% conversion: kordoc 4.24.51 -> tools/gen1.py (auto); review=proofread-sonnet; flags=none
\begin{question}{33}
법률행위의 무효와 취소에 관한 설명으로 옳지 \uline{않은} 것은? (다툼이 있으면 판례에 의함)
\begin{choices}
  \item 법률행위의 무효는 법률행위의 성립을 전제로 한다.
  \item 무효인 가등기를 유효한 것으로 전용하기로 약정하였다면 가등기는 소급하여 유효한 등기로 된다.
  \item 당사자가 입양의사로 친생자 출생신고를 한 경우 입양의 실질적 요건이 모두 구비되었다면 입양의 효력이 발생한다.
  \item 강박상태에서 벗어나지 않은 상태에서 한 추인은 그 효력이 발생하지 않는다.
  \item 신용카드 이용계약 취소시 미성년자가 반환할 금전상 이익은 특별한 사정이 없는 한 현존하는 것으로 추정된다.
\end{choices}
\end{question}
\answer{2}
